\documentclass[letterpaper]{article}
\usepackage[utf8]{inputenc}
\usepackage{setspace}
\usepackage{latexsym,amsfonts}
\usepackage{soul, color}
\usepackage{multicol}
\usepackage[colorlinks=true,linkcolor=blue]{hyperref}
\sethlcolor{yellow}
\title{\textbf{MSC2020-Mathematics Subject Classification System}}
\author{\textbf{Associate Editors of Mathematical Reviews and zbMATH}}
\date{}
\voffset -1in
\hoffset -1in
\textwidth 7.0in
\textheight 9in
\begin{document}
\maketitle
\small
\begin{multicols}{2}
\begin{description}
\item []\hyperref[00-XX]{\textbf{00}}  General and overarching topics; collections 
\item []\hyperref[01-XX]{\textbf{01}}  History and biography  
\item []\hyperref[03-XX]{\textbf{03}}  Mathematical logic and foundations 
\item []\hyperref[05-XX]{\textbf{05}}  Combinatorics  
\item []\hyperref[06-XX]{\textbf{06}}  Order, lattices, ordered algebraic 
structures 
\item []\hyperref[08-XX]{\textbf{08}}  General algebraic systems 
\item []\hyperref[11-XX]{\textbf{11}}  Number theory 
\item []\hyperref[12-XX]{\textbf{12}}  Field theory and polynomials 
\item []\hyperref[13-XX]{\textbf{13}}  Commutative algebra 
\item []\hyperref[14-XX]{\textbf{14}}  Algebraic geometry 
\item []\hyperref[15-XX]{\textbf{15}}  Linear and multilinear algebra; matrix 
theory 
\item []\hyperref[16-XX]{\textbf{16}}  Associative rings and algebras  
\item []\hyperref[17-XX]{\textbf{17}}  Nonassociative rings and algebras 
\item []\hyperref[18-XX]{\textbf{18}}  Category theory; homological algebra 
\item []\hyperref[19-XX]{\textbf{19}}  $K$-theory  
\item []\hyperref[20-XX]{\textbf{20}}  Group theory and generalizations 
\item []\hyperref[22-XX]{\textbf{22}}  Topological groups, Lie groups  
\item []\hyperref[26-XX]{\textbf{26}}  Real functions 
\item []\hyperref[28-XX]{\textbf{28}}  Measure and integration  
\item []\hyperref[30-XX]{\textbf{30}}  Functions of a complex variable  
\item []\hyperref[31-XX]{\textbf{31}}  Potential theory  
\item []\hyperref[32-XX]{\textbf{32}}  Several complex variables and analytic 
spaces  
\item []\hyperref[33-XX]{\textbf{33}}  Special functions 
\item []\hyperref[34-XX]{\textbf{34}}  Ordinary differential equations 
\item []\hyperref[35-XX]{\textbf{35}}  Partial differential equations 
\item []\hyperref[37-XX]{\textbf{37}}  Dynamical systems and ergodic theory  
\item []\hyperref[39-XX]{\textbf{39}}  Difference and functional equations 
\item []\hyperref[40-XX]{\textbf{40}}  Sequences, series, summability 
\item []\hyperref[41-XX]{\textbf{41}}  Approximations and expansions  
\item []\hyperref[42-XX]{\textbf{42}}  Harmonic analysis on Euclidean spaces 
\item []\hyperref[43-XX]{\textbf{43}}  Abstract harmonic analysis  
\item []\hyperref[44-XX]{\textbf{44}}  Integral transforms, operational calculus 
\item []\hyperref[45-XX]{\textbf{45}}  Integral equations 
\item []\hyperref[46-XX]{\textbf{46}}  Functional analysis 
\item []\hyperref[47-XX]{\textbf{47}}  Operator theory 
\item []\hyperref[49-XX]{\textbf{49}}  Calculus of variations and optimal 
control; 
optimization
\item []\hyperref[51-XX]{\textbf{51}}  Geometry  
\item []\hyperref[52-XX]{\textbf{52}}  Convex and discrete geometry 
\item []\hyperref[53-XX]{\textbf{53}}  Differential geometry  
\item []\hyperref[54-XX]{\textbf{54}}  General topology  
\item []\hyperref[55-XX]{\textbf{55}}  Algebraic topology 
\item []\hyperref[57-XX]{\textbf{57}}  Manifolds and cell complexes  
\item []\hyperref[58-XX]{\textbf{58}}  Global analysis, analysis on manifolds 
\item []\hyperref[60-XX]{\textbf{60}}  Probability theory and stochastic 
processes  
\item []\hyperref[62-XX]{\textbf{62}}  Statistics 
\item []\hyperref[65-XX]{\textbf{65}}  Numerical analysis 
\item []\hyperref[68-XX]{\textbf{68}}  Computer science  
\item []\hyperref[70-XX]{\textbf{70}}  Mechanics of particles and systems  
\item []\hyperref[74-XX]{\textbf{74}}  Mechanics of deformable solids 
\item []\hyperref[76-XX]{\textbf{76}}  Fluid mechanics  
\item []\hyperref[78-XX]{\textbf{78}}  Optics, electromagnetic theory  
\item []\hyperref[80-XX]{\textbf{80}}  Classical thermodynamics, heat transfer  
\item []\hyperref[81-XX]{\textbf{81}}  Quantum theory 
\item []\hyperref[82-XX]{\textbf{82}}  Statistical mechanics, structure of 
matter 
\item []\hyperref[83-XX]{\textbf{83}}  Relativity and gravitational theory 
\item []\hyperref[85-XX]{\textbf{85}}  Astronomy and astrophysics 
\item []\hyperref[86-XX]{\textbf{86}}  Geophysics 
\item []\hyperref[90-XX]{\textbf{90}}  Operations research, mathematical 
programming 
\item []\hyperref[91-XX]{\textbf{91}}  Game theory, economics, social and 
behavioral 
sciences
\item []\hyperref[92-XX]{\textbf{92}}  Biology and other natural sciences 
\item []\hyperref[93-XX]{\textbf{93}}  Systems theory; control  
\item []\hyperref[94-XX]{\textbf{94}}  Information and communication, circuits 
\item []\hyperref[97-XX]{\textbf{97}}  Mathematics education 
\end{description}
\end{multicols}
\newpage
\normalsize
\begin{doublespace}
This  document  is  a  printed  form  of  MSC2020,  an  MSC  revision  produced
jointly by the editorial staffs of Mathematical Reviews (MR) and Zentralblatt 
f\"{u}r Mathematik (zbMATH) in consultation with the mathematical community. 
The goals of
this revision of the Mathematics Subject Classification (MSC) were set out in 
the
announcement of it and call for comments by the Executive Editor of MR and the
Chief Editor of zbMATH in July 2016. This document results from the MSC 
revision
process that has been going on since then. MSC2020 will be fully deployed from
January 2020.
\par
The editors of MR and zbMATH deploying this revision therefore ask for feedback 
on
remaining errors to help in this work, which should be given through e-mail to
\href{mailto:feedback@msc2020.org}{feedback@msc2020.org}. They are grateful for the many suggestions that were received previously, which have greatly influenced what we have.
\newpage
\begin{center}\textbf{\Large{How to use the
Mathematics Subject Classification [MSC]}}\end{center}
The main purpose of the classification of items in the mathematical literature
using  the  Mathematics  Subject  Classification  scheme  is  to  help  users  
find  the
items of present or potential interest to them as readily as possible—in 
products
derived from the Mathematical Reviews Database (MRDB) such as MathSciNet, in
Zentralblatt MATH (zbMATH), or anywhere else where this classification scheme is
used. An item in the mathematical literature should be classified so as to 
attract the
attention of all those possibly interested in it. The item may be something 
that
falls  squarely  within  one  clear  area  of  the  MSC,  or  it  may  involve  
several  areas.
Ideally, the MSC codes attached to an item should represent the subjects to 
which
the item contains a contribution. The classification should serve both those 
closely
concerned with specific subject areas, and those familiar enough with subjects 
to
apply their results and methods elsewhere, inside or outside of mathematics. It 
will
be extremely useful for both users and classifiers to familiarize themselves 
with the
entire  classification  system  and  thus  to  become  aware  of  all  the  
classifications  of
possible interest to them.
Every   item   in   the   MRDB   or   zbMATH   receives   precisely   one
primary
classification,   which   is   simply   the   MSC   code   that   describes   
its   principal contribution.  When  an  item  contains  several  principal  contributions  to  
different
areas, the primary classification should cover the most important among them. A
paper or book may be assigned one or several secondary classification numbers to
cover any remaining principal contributions, ancillary results, motivation or 
origin of
the matters discussed, intended or potential field of application, or other 
significant
aspects worthy of notice.
The principal contribution is meant to be the one including the most important
part of the work actually done in the item. For example, a paper whose main 
overall
content  is  the  solution  of  a  problem  in  graph  theory,  which  arose  in 
 computer
science  and  whose  solution  is  (perhaps)  at  present  only  of  interest  
to  computer
scientists,  would  have  a  primary  classification  in  05C  (Graph  Theory)  
with  one
or  more  secondary  classifications  in  68  (Computer  Science);  conversely,  
a  paper
whose  overall  content  lies  mainly  in  computer  science  should  receive  a 
 primary
classification in 68, even if it makes heavy use of graph theory and proves 
several
new graph-theoretic results along the way.
There are two types of cross-references given at the end of many of the MSC2020 entries
in the MSC. The first type is in braces: “\{For A, see X\}”; if this appears in section
Y,  it  means  that  contributions  described  by  A  should  usually  be  
assigned  the
classification  code  X,  not  Y.  The  other  type  of  cross-reference  
merely 
 points  out
related  classifications;  it  is  in  brackets:  “[See  also
...
]”,  “[See  mainly
...
]”,  etc.,
and the classification codes listed in the brackets may, but need not, be 
included in
the classification codes of a paper, or they may be used in place of the 
classification
where the cross-reference is given. The classifier must judge which 
classification is
the most appropriate for the paper at hand.
\end{doublespace}
\newpage
%\begin{multicols}{2}
